\documentclass[aspectratio=169,12pt]{beamer}

% --- Theme ---
\usetheme{metropolis}
\usepackage{appendixnumberbeamer}
\usepackage{booktabs}
\usepackage{graphicx}
\usepackage{xcolor}
\usepackage{tikz}
\usepackage{amsmath}

% --- ICLR color scheme ---
\definecolor{iclrprimary}{HTML}{7B2D8E}
\definecolor{iclraccent}{HTML}{F0AD00}
\setbeamercolor{frametitle}{bg=iclrprimary, fg=white}
\setbeamercolor{progress bar}{fg=iclraccent}
\setbeamercolor{alerted text}{fg=iclraccent}

% --- Metadata ---
\title{Who's Afraid of Communist AI:\\Epistemic Transfer and Ideological Convergence in LMs}
\subtitle{ICLR 2026}
\author{Anonymous ICLR Submission}
\institute{}
\date{}

\setbeameroption{hide notes}

\begin{document}

% ============================================================
% SLIDE 1: Title
% ============================================================
\maketitle

\note{
  [30 sec] Welcome. Today I'll present our work on how SFT transfers not just vocabulary but deep epistemic priors in language models.
}

% ============================================================
% SLIDE 2: Problem
% ============================================================
\begin{frame}{Problem: What Does SFT Really Transfer?}
  \begin{itemize}
    \item SFT is the \textbf{foundational} step in LLM alignment
    \item Alignment literature focuses on safety, helpfulness, preference consistency
    \item \alert{Open question:} Does SFT on ideologically structured data transfer \textit{reasoning patterns} beyond the training domain?
    \item Prior work documents ideological priors as a \textit{static} property \cite{kronlund2024propaganda,lee2026ideological}
  \end{itemize}

  \vspace{0.5em}
  \begin{block}{Our Dynamic Question}
    Can targeted SFT on a \textit{non-dominant} analytical framework measurably shift reasoning --- and what \textbf{collateral effects} does this produce?
  \end{block}

  \note{
    [45 sec] SFT is how we align models, but we don't understand what it transfers beyond the training domain. Prior work showed LLMs carry ideological priors from pretraining, but treated it as static. We asked: if we fine-tune on a specific analytical framework, does it shift reasoning in unseen domains? What are the collateral effects?
    Transition: "To answer this, we fine-tuned Qwen3.5-9B on Dialectical Materialist analysis."
  }
\end{frame}

% ============================================================
% SLIDE 3: Methodology
% ============================================================
\begin{frame}{Methodology: DM-Aligned SFT}
  \begin{columns}[T]
    \begin{column}{0.48\textwidth}
      \textbf{Training Setup}
      \begin{itemize}
        \item Student: Qwen3.5-9B (Instruct)
        \item Teacher: Qwen3.5-27B (data gen. only)
        \item QLoRA NF4, $r\!=\!32$, 7 target modules
        \item 1,460 DM Q\&A pairs, 3 epochs
        \item Single RTX 5090 (32GB)
      \end{itemize}
    \end{column}
    \begin{column}{0.48\textwidth}
      \textbf{Evaluation} (bf16, 0-shot greedy)
      \begin{itemize}
        \item \textit{General:} MMLU, HumanEval, GPQA
        \item \textit{Formal causal:} Corr2Cause (1,162 samples)
        \item \textit{Applied causal:} EconCausal (3,943 samples)
      \end{itemize}
    \end{column}
  \end{columns}

  \note{
    [45 sec] We fine-tuned Qwen3.5-9B via QLoRA on 1,460 question-answer pairs generated by a 27B teacher using Dialectical Materialist analysis prompts. We evaluated across three domains: general capability, formal causal logic, and applied economic causal reasoning. All in bf16 with 0-shot greedy decoding.
    Transition: "Here's what we found --- and it's surprising."
  }
\end{frame}

% ============================================================
% SLIDE 4: Key Finding --- Three Divergent Outcomes
% ============================================================
\begin{frame}{Three Divergent Outcomes}
  \begin{columns}[T]
    \begin{column}{0.32\textwidth}
      \begin{block}{General Capability}
        \centering\Large\textbf{Preserved}
      \end{block}
      \vspace{0.3em}
      \begin{itemize}
        \item MMLU: $-0.8$pp
        \item HumanEval: $0.0$pp
        \item GPQA: $-1.5$pp
        \item All within binomial variance
      \end{itemize}
    \end{column}
    \begin{column}{0.32\textwidth}
      \begin{block}{Formal Causal Reasoning}
        \centering\Large\textbf{+38.3pp}
      \end{block}
      \vspace{0.3em}
      \begin{itemize}
        \item Corr2Cause
        \item Corrects 520 baseline errors
        \item Introduces only 75 new errors
        \item Net gain: \textbf{+445}
      \end{itemize}
    \end{column}
    \begin{column}{0.32\textwidth}
      \begin{block}{Applied Causal Reasoning}
        \centering\Large\textbf{-12.4pp}
      \end{block}
      \vspace{0.3em}
      \begin{itemize}
        \item EconCausal Task1: $\mathbf{-12.4}$pp
        \item Task1 Finance: $\mathbf{-13.5}$pp
        \item Task3: $\mathbf{-10.8}$pp
        \item All highly significant
      \end{itemize}
    \end{column}
  \end{columns}

  \note{
    [60 sec] Three outcomes. First: general capability is preserved. MMLU drops 0.8 points, HumanEval is identical, GPQA drops 1.5 points --- all within statistical noise. Second: formal causal reasoning improves dramatically. Corr2Cause jumps 38.3 points, with a net gain of 445 correct answers. Third: applied economic causal reasoning regresses severely --- 12 to 13.5 points across tasks, all highly significant.
    Transition: "How can the same training improve one causal benchmark while breaking another? The answer is a single emergent artifact."
  }
\end{frame}

% ============================================================
% SLIDE 5: The Hedging Artifact
% ============================================================
\begin{frame}{The Hedging Artifact}
  \begin{alertblock}{Dominant Failure Mode: Positive-to-Mixed Hedging}
    The finetuned model converts correct \textbf{``+''} predictions to ambiguous \textbf{``mixed''} answers.
  \end{alertblock}

  \vspace{0.5em}
  \begin{columns}[T]
    \begin{column}{0.55\textwidth}
      \begin{itemize}
        \item Task1 Econ: $+$ to mixed = 96 of 182 regressions (\textbf{52.7\%})
        \item Task1 Finance: 95 of 174 (\textbf{54.6\%})
        \item Positive-effect conversion: \textbf{77--85\%} of all Task1 regressions
        \item \alert{Hedging absent from training data} (teacher rate: 4.0\%)
      \end{itemize}
    \end{column}
    \begin{column}{0.4\textwidth}
      \begin{block}{Mechanism}
        \textit{Transferred Epistemic Prior}
        \vspace{0.3em}

        Model learns DM principles:
        \begin{itemize}
          \item ``Outcomes depend on material conditions''
          \item ``Effects mediated by power relations''
          \item ``Cannot reduce to single direction''
        \end{itemize}
        Applies skepticism \textit{indiscriminately}.
      \end{block}
    \end{column}
  \end{columns}

  \note{
    [60 sec] The dominant regression pattern is positive-to-mixed hedging. The model takes correct positive causal predictions and converts them to ambiguous mixed answers. On Task1 Econ, this accounts for 52.7% of all regressions. Here's the key: this hedging is ABSENT from the training data. The teacher's hedging rate is only 4%. It's an emergent artifact. The model learned DM's structural skepticism principles and applies them universally --- even to empirical economics questions where definitive directional effects are the norm. This is not a reasoning error. It's a transferred epistemic prior.
    Transition: "This explains why Corr2Cause improves while EconCausal regresses."
  }
\end{frame}

% ============================================================
% SLIDE 6: Why the Divergence?
% ============================================================
\begin{frame}{Why Corr2Cause Improves While EconCausal Regresses}
  \begin{columns}[T]
    \begin{column}{0.48\textwidth}
      \begin{block}{Corr2Cause: \textbf{+38.3pp}}
        \textit{Formal structural analysis}
      \end{block}
      \vspace{0.3em}
      \begin{itemize}
        \item Conditional independence reasoning
        \item DM training \textit{strengthens}: trace dependencies, identify confounders, evaluate causal pathways
        \item Non-child descendant: $13.0\% \to 94.3\%$ (\textbf{+81.3pp})
        \item Non-parent ancestor: $14.4\% \to 87.2\%$ (\textbf{+72.8pp})
      \end{itemize}
    \end{column}
    \begin{column}{0.48\textwidth}
      \begin{block}{EconCausal: \textbf{-12.4pp}}
        \textit{Applied causal identification}
      \end{block}
      \vspace{0.3em}
      \begin{itemize}
        \item Predicting directional effects from empirical context
        \item DM skepticism prior \textit{overrides} empirical claims
        \item Model learns ``everything depends on structural conditions''
        \item Applies to questions where correct answer is a definitive $+$ or $-$
      \end{itemize}
    \end{column}
  \end{columns}

  \vspace{0.5em}
  \begin{block}{Confirmation}
    Mirrors \citet{syrgkanis2026causalreasoning}: causal reasoning has \textbf{two distinct bottlenecks} --- formal logic (strengthened) and identification details (weakened).
  \end{block}

  \note{
    [45 sec] Corr2Cause requires formal structural analysis --- reasoning about conditional independence given a graph. DM training strengthens this: the model learns to trace dependencies, identify confounders, evaluate causal pathways. Gains of 72 to 81 points on complex templates. EconCausal requires applied identification --- predicting directional effects from empirical context. Here, DM's skepticism prior overrides empirical claims. The model learned ``everything depends on structural conditions'' and applies it where the correct answer is a definitive plus or minus. This confirms Syrgkanis's finding: causal reasoning has two distinct bottlenecks.
    Transition: "What does this mean for alignment research?"
  }
\end{frame}

% ============================================================
% SLIDE 7: Implications
% ============================================================
\begin{frame}{Implications for Alignment Research}
  \begin{enumerate}
    \item \textbf{SFT transfers epistemic priors, not just behavior}
    \begin{itemize}
      \item Hedging artifact is absent from training data, emerges from internalized reasoning framework
      \item Shift operates at level of reasoning heuristics, not response patterns
    \end{itemize}

    \item \textbf{Standard alignment may carry unaudited priors}
    \begin{itemize}
      \item RLHF emphasizing helpfulness/harmlessness may transfer deference to authority, controversy avoidance
      \item These go unaudited because they align with evaluator's own priors
    \end{itemize}

    \item \textbf{The Identification-Logic Split as a diagnostic}
    \begin{itemize}
      \item If an intervention improves formal logic but impairs applied identification $\to$ over-generalized structural skepticism
      \item Benchmarks need \textit{both} formal and applied causal tasks
    \end{itemize}
  \end{enumerate}

  \note{
    [45 sec] Three implications. First: SFT transfers epistemic priors, not just behavioral patterns. The shift operates at the level of reasoning heuristics. Second: if DM training produces detectable epistemic transfer, then standard RLHF likely carries similar hidden priors --- deference to authority, controversy avoidance --- that go unaudited because they align with the evaluator's assumptions. Third: the Corr2Cause/EconCausal divergence gives us a diagnostic tool. If an intervention improves formal logic while impairing applied identification, it may be over-generalizing structural skepticism.
    Transition: "Let me summarize."
  }
\end{frame}

% ============================================================
% SLIDE 8: Summary
% ============================================================
\begin{frame}{Summary}
  \begin{enumerate}
    \item \textbf{Problem:} SFT's effects on reasoning beyond training domain are poorly understood
    \item \textbf{Method:} QLoRA SFT on 1,460 DM-aligned Q\&A pairs (Qwen3.5-9B)
    \item \textbf{Finding 1:} General capability preserved (MMLU $-0.8$pp, HumanEval $0.0$pp)
    \item \textbf{Finding 2:} Formal causal reasoning improves dramatically (Corr2Cause $+38.3$pp)
    \item \textbf{Finding 3:} Applied causal reasoning regresses (EconCausal $-12.4$pp)
    \item \textbf{Mechanism:} Emergent hedging bias --- transferred epistemic prior, absent from training data
    \item \textbf{Implication:} Alignment training transfers reasoning priors that should be explicitly audited
  \end{enumerate}

  \note{
    [30 sec] We show that SFT on ideologically structured data produces measurable, domain-specific reasoning shifts without catastrophic general degradation. The key finding: alignment training transfers epistemic priors. The finetuned model learned structural skepticism and applied it universally, improving formal logic while impairing applied identification. If DM training produces detectable transfer, standard alignment likely carries unaudited priors we don't yet measure.
  }
\end{frame}

% ============================================================
% SLIDE 9: Future Work
% ============================================================
\begin{frame}{Future Work: Breaking the Hedging Prior}
  \begin{itemize}
    \item \textbf{GRPO with process rewards} (RLVMR framework \cite{zhang2025rlvmr})
    \item Dual advantage: trajectory-level outcome advantage + meta-reasoning advantage
    \item Reward definitive commitment when warranted
    \item Goal: break hedging prior while maintaining structural reasoning quality
  \end{itemize}

  \vspace{0.5em}
  \begin{block}{Current Status}
    \begin{itemize}
      \item GRPO v3 (outcome-only): stalled at 806/1500 steps, no reward improvement
      \item GRPO v4 (process rewards): tagged reasoning with rule-based rewards for planning, commitment, reflection, monitoring
    \end{itemize}
  \end{block}

  \note{
    [30 sec] Our ongoing work explores GRPO with process rewards to break the hedging prior. We use dual advantage: outcome correctness plus meta-reasoning rewards for definitive commitment. The outcome-only track stalled, suggesting reward signal alone is insufficient. The process reward track explicitly rewards commitment when warranted, while maintaining structural reasoning quality.
    Transition: "Thank you. I'd be happy to take questions."
  }
\end{frame}

% ============================================================
% SLIDE 10: Thank You
% ============================================================
\begin{frame}{Thank You}
  \begin{center}
    \vspace{1.5em}
    \Large Questions? \\[2em]
    \textbf{Key Takeaway:} \\
    SFT transfers epistemic priors, not just behavior. \\
    Alignment training carries hidden reasoning shifts that should be audited. \\[2em]
    Code: \url{github.com/...}
  \end{center}

  \note{
    Leave this slide up during Q\&A. Have backup slides ready for:
    - Full MMLU subtask breakdown
    - Corr2Cause template-level analysis
    - EconCausal sample-level regression patterns
    - Training hyperparameters
    - Dataset quality metrics
  }
\end{frame}

% ============================================================
% BACKUP SLIDES
% ============================================================
\appendix

\begin{frame}{Backup: Full General Capability Results}
  \begin{center}
  \begin{tabular}{lccc}
  \toprule
  Benchmark & Baseline & Finetuned & $\Delta$ \\
  \midrule
  MMLU Overall & $78.72\%$ & $77.96\%$ & $-0.76$pp \\
  MMLU STEM & $78.53\%$ & $78.24\%$ & $-0.29$pp \\
  MMLU Social Sci. & $86.68\%$ & $86.19\%$ & $-0.49$pp \\
  MMLU Humanities & $70.69\%$ & $69.86\%$ & $-0.83$pp \\
  HumanEval pass@1 & $70.73\%$ & $70.73\%$ & $0.00$pp \\
  GPQA Diamond & $47.47\%$ & $45.96\%$ & $-1.51$pp \\
  IFEval (strict) & $45.84\%$ & $44.64\%$ & $-1.20$pp \\
  IFEval (loose) & $49.35\%$ & $47.60\%$ & $-1.75$pp \\
  \bottomrule
  \end{tabular}
  \end{center}
  \note{Use if asked about general capability preservation. All deltas within one standard error of zero.}
\end{frame}

\begin{frame}{Backup: Corr2Cause Template-Level Gains}
  \begin{center}
  \begin{tabular}{lccccc}
  \toprule
  Template & GT True\% & Baseline True\% & Baseline Acc & Finetuned Acc & $\Delta$ \\
  \midrule
  child & $0.0\%$ & $19.1\%$ & $80.9\%$ & $100.0\%$ & $+19.1$pp \\
  non-child desc. & $5.7\%$ & $92.7\%$ & $13.0\%$ & $94.3\%$ & $+81.3$pp \\
  non-parent anc. & $11.3\%$ & $96.9\%$ & $14.4\%$ & $87.2\%$ & $+72.8$pp \\
  has\_confounder & $8.8\%$ & $97.4\%$ & $11.9\%$ & $54.9\%$ & $+43.0$pp \\
  parent & $33.5\%$ & $43.8\%$ & $62.9\%$ & $66.5\%$ & $+3.6$pp \\
  has\_collider & $33.1\%$ & $98.4\%$ & $34.7\%$ & $44.6\%$ & $+9.8$pp \\
  \bottomrule
  \end{tabular}
  \end{center}
  \note{Use if asked about baseline True-bias or per-template breakdown.}
\end{frame}

\begin{frame}{Backup: EconCausal Regression Patterns}
  \begin{center}
  \begin{tabular}{lccccl}
  \toprule
  Pattern & T1 Econ & T1 Fin & T2 & T3 & Description \\
  \midrule
  $+ \to$ mixed & $96$ & $95$ & $9$ & $34$ & Hedging positive effects \\
  $+ \to -$ & $37$ & $41$ & $5$ & $36$ & Flipping positive to negative \\
  $+ \to$ none & $25$ & $20$ & $-$ & $-$ & Nullifying positive effects \\
  $- \to$ mixed & $8$ & $-$ & $-$ & $11$ & Hedging negative effects \\
  $- \to +$ & $6$ & $-$ & $-$ & $-$ & Flipping negative to positive \\
  \bottomrule
  \end{tabular}
  \end{center}
  \note{Use if asked about specific regression patterns. Positive-to-mixed dominates all tasks.}
\end{frame}

\begin{frame}{Backup: Training Configuration}
  \begin{center}
  \begin{tabular}{ll}
  \toprule
  Parameter & Value \\
  \midrule
  Base model & Qwen3.5-9B (Instruct) \\
  Quantization & NF4 (bitsandbytes) \\
  LoRA rank ($r$) & $32$ \\
  LoRA $\alpha$ & $32$ \\
  LoRA dropout & $0.05$ \\
  Target modules & $7$ \\
  Learning rate & $2\times10^{-4}$ \\
  Scheduler & Cosine (warmup 100 steps) \\
  Max steps & $1,000$ \\
  Epochs & $3$ \\
  Batch size & $1$ (effective $4$) \\
  Neftune noise $\alpha$ & $5$ \\
  Hardware & RTX 5090 (32GB VRAM) \\
  \bottomrule
  \end{tabular}
  \end{center}
  \note{Use if asked about training details.}
\end{frame}

\end{document}
